\section{Graph alphabets and the coloring reduction}
\label{graph:section}

This section reduces the elementary expansion to coefficients of
monomial symmetric functions evaluated in a quantum torus. It also
fixes the permutation convention used by the transport construction.

\subsection{Admissible words and insertion digits}

For the graph $G$ in the introduction put
$l_a=\#\{b<a:\{b,a\}\in E(G)\}$ and $M=\sum_a l_a$.
The earlier neighbors of $a$ are a suffix of $[a-1]$, and they form a
clique. Let
\[
 \mathcal A_G=\{w\in S_n:
    w_i<w_{i+1}\Longrightarrow\{w_i,w_{i+1}\}\in E(G)\}.
\]
We call these words \emph{admissible}. Their forbidden adjacent pairs
are nonedge ascents. Write $\inv_G(w)$ for their edge-inversion count,
given by the same formula as $\ginv_G$.

\begin{lemma}[Insertion digits]
\label{graph:digits}
There is a bijection between $\mathcal A_G$ and the sequences
$(j_1,\ldots,j_n)$ with $0\le j_a\le l_a$. Under this bijection,
$\inv_G(w)=\sum_a j_a$. Reversal is a bijection
$D_G^0\to\mathcal A_G$, and
\begin{equation}
 \inv_G(\operatorname{reverse}\sigma)=M-\ginv_G(\sigma).
 \label{graph:reverse}
\end{equation}
\end{lemma}

\begin{proof}
Delete the largest letter $a$ from an admissible word on $[a]$.
If it has predecessor $b$ and successor $c$, then $b<a$ and
$h(b)\ge a>c$, so the newly consecutive pair $(b,c)$ is allowed.
The cases without a predecessor or successor are immediate. Conversely,
$a$ can be inserted at the beginning or immediately after any earlier
neighbor, and these are exactly the allowed positions. At such a
position the new inversions are the earlier neighbors lying to its
right. Reading those neighbors in their order in the shorter word
shows that the available counts are $l_a,l_a-1,\ldots,0$, each once.
Successive insertion and deletion prove the bijection and degree
formula. Reversal exchanges nonedge descents and nonedge ascents and
reverses the order of the endpoints of every edge, proving
Equation~\eqref{graph:reverse}.
\end{proof}

The relation $a<b$ and $\{a,b\}\notin E(G)$ is transitive: it is
equivalent to $h(a)<b$. Consequently, an increasing word in which each
adjacent pair is a nonedge lists an independent set. This observation
will identify packets of independent vertices with required ascent
positions in Section~\ref{route:section}.

\subsection{The quantum torus of an ordered graph}

For this subsection let $\Gamma$ be any natural unit interval graph on
a finite ordered vertex set. Let $L_\Gamma$ be the free abelian group
on its vertices, identifying a vertex with its basis vector. Define an
integer alternating form by
\begin{equation}
 \Omega(b,a)=
 \begin{cases}1,&a<b\text{ and }\{a,b\}\in E(\Gamma),\\
                0,&a<b\text{ and }\{a,b\}\notin E(\Gamma),
 \end{cases}
 \label{graph:form}
\end{equation}
and extend by bilinearity and skew symmetry. Over
$\Q((v^{-1}))$, the quantum torus has basis $X_m$, $m\in L_\Gamma$,
and product
\begin{equation}
 X_mX_{m'}=v^{\Omega(m,m')}X_{m+m'}.
 \label{graph:torus}
\end{equation}
The coefficient field consists of Laurent series with bounded upper
powers of $v$. All elements considered in this section are finite
Laurent polynomials; the additional cone completions will be specified
when needed.

For a set $I$ of vertices write $a_I=\sum_{a\in I}a$, and define
\begin{equation}
 E_k(\Gamma)=\sum_{\substack{I\subseteq V(\Gamma)\text{ independent}\\|I|=k}}
 X_{a_I},\qquad E_0(\Gamma)=1.
 \label{graph:elementaries}
\end{equation}

The graph alphabet used here is a centered quantum-torus version of
noncommutative poset symmetric functions \cite{Hwa24,BEPS25}.  More
precisely, on $V(\Gamma)$ define $a<_P b$ when $a<b$ and $\{a,b\}$ is a nonedge.
The chain elementary sums in \cite[Equation~(2.4)]{BEPS25} become our
independent-set sums, and the relations in
\cite[Definition~3.16]{BEPS25} become the relations above under
$u_a\mapsto X_a$ and $t\mapsto v^2$.  We prove commutation directly
in these conventions.

\begin{proposition}[Commuting independent-set sums]
\label{graph:commutation}
For every natural unit interval graph $\Gamma$ and all $j,k\ge0$,
$E_j(\Gamma)E_k(\Gamma)=E_k(\Gamma)E_j(\Gamma)$.
\end{proposition}

\begin{proof}
We adapt the two-color involution underlying
\cite[Lemma~4.4 and Theorem~4.5]{SW16}.
A term on the left is an ordered pair of independent sets $(I,J)$,
with exponent $a_I+a_J$ and weight $v^{\Omega(a_I,a_J)}$.
Vertices in $I\cap J$ have no neighbors in $I\cup J$ and contribute
no skew weight. In the remaining induced graph, every earlier-neighbor
set and every later-neighbor set is a clique. Since the graph is
bipartite with parts $I\setminus J$ and $J\setminus I$, each vertex
has at most one earlier and one later neighbor. More explicitly, an
edge that skipped a vertex of $I\cup J$ in the total order would make
a triangle: the interval-closure property supplies both intervening
edges. Thus each nontrivial component is an increasing path with
alternating membership in $I$ and $J$.

On every component with an odd number of vertices, exchange the two
memberships; keep the even components and the doubled isolated
vertices unchanged. Every even component has equally many vertices in
the two sets. Every odd component reverses its size difference, so the
new ordered pair $(I',J')$ has sizes $k,j$. On an odd component the
successive edge contributions to $\Omega(a_I,a_J)$ alternate between
$1$ and $-1$, with equally many of each. Its skew contribution is
therefore zero before and after the exchange. There are no edges
between components. Exponents and weights are preserved, and repeating
the exchanges recovers $(I,J)$. This is the required bijection of
terms.
\end{proof}

Let $\alpha(\Gamma)$ denote the maximum size of an independent set in
$\Gamma$, and suppose $\alpha(\Gamma)\le D$. The ring of symmetric polynomials in
$D$ commuting variables is the polynomial ring in $e_1,\ldots,e_D$.
Proposition~\ref{graph:commutation} therefore defines an algebra
homomorphism by $e_k\mapsto E_k(\Gamma)$, with $E_k=0$ for $k>D$.
We denote its value on a symmetric polynomial $F$ by $F(Y)$, and
write $e_k(Y)=E_k$. The symbol $Y$ denotes this evaluation; no
individual roots $y_i$ in the quantum torus are being asserted.
In particular, $m_\lambda(Y)$ is well defined for partitions of length
at most $D$, where $m_\lambda$ is the sum of distinct monomials with
exponent multiset $\lambda$.

\subsection{The elementary coefficient identity}

Let $G$ be the induced ordered graph on distinct vertices
$a_1,\ldots,a_n$ of $\Gamma$, and put $u=a_1+\cdots+a_n$.
For a proper coloring $f$ define
$\operatorname{inv}_G(f)=\#\{a<b:\{a,b\}\in E(G),\ f(a)>f(b)\}$
and let $X_G(X;q)$ be the corresponding proper-color inversion
generating function.

The following coefficient reduction is the elementary/monomial form
of the noncommutative Cauchy pairing in
\cite[Theorem~4.4, Corollary~4.5 and Proposition~4.25]{Hwa24} and
\cite[Proposition~2.9 and Corollary~2.10]{BEPS25}.  We include its
proof to fix the graph-inversion and centered-monomial factors.

\begin{proposition}[Coloring kernel]
\label{graph:kernel}
If $\alpha(\Gamma)\le D$, then $X_G(X;q)$ is symmetric and
\begin{equation}
 X_G(X;v^2)=\sum_{\substack{\lambda\vdash n\\\ell(\lambda)\le D}}
       v^M[X_u]m_\lambda(Y)\,e_\lambda(X).
 \label{graph:coefficient}
\end{equation}
The equality is independent of the extra vertices of $\Gamma$.
\end{proposition}

\begin{proof}
Work first with $r$ colors. Expand the product, in increasing color
order,
\[
 \prod_{c=1}^r\left(\sum_{k=0}^D x_c^k E_k\right).
\]
Its $X_u$ coefficient chooses a disjoint independent set for each
color, covering precisely the original vertices. Extra vertices
cannot occur because each exponent in each factor has nonnegative
vertex coordinates. These choices are exactly the proper colorings.
An edge whose later vertex has smaller color contributes $+1$ to the
torus exponent, and an edge whose earlier vertex has smaller color
contributes $-1$. The total exponent is
$2\operatorname{inv}_G(f)-M$. Hence multiplying this coefficient by
$v^M$ gives $X_G(x_1,\ldots,x_r;v^2)$.

For ordinary variables $y_1,\ldots,y_D$ the elementary dual kernel
identity is
\[
 \prod_{c=1}^r\prod_{i=1}^D(1+x_cy_i)
   =\sum_\lambda e_\lambda(x_1,\ldots,x_r)m_\lambda(y_1,\ldots,y_D).
\]
Indeed, select for each labeled $y_i$ a subset of the $x$ variables;
grouping its subset sizes by their multiset gives the displayed sum.
Evaluate the symmetric $y$ polynomials at $Y$, extract $X_u$, and
multiply by $v^M$. This proves Equation~\eqref{graph:coefficient} for
every $r$, and therefore as a symmetric-function identity. The
coloring interpretation proves independence of the ambient graph.
\end{proof}

\begin{example}[The centered normalization on two vertices]
Let $\Gamma=G$ have ordered vertices $a<b$. If they are adjacent,
$E_2=0$ and $m_{(2)}(Y)=E_1^2$ has $X_{a+b}$-coefficient
$v+v^{-1}$. The factor $v^M=v$ in
Equation~\eqref{graph:coefficient} therefore gives
$X_G(X;q)=(1+q)e_2(X)$.
If they are nonadjacent, $E_2=X_{a+b}$ and the two terms in
$E_1^2$ at this exponent each have weight one. Hence
$[X_{a+b}]m_{(2)}(Y)=[X_{a+b}](E_1^2-2E_2)=0$, whereas
$[X_{a+b}]m_{(1,1)}(Y)=[X_{a+b}]E_2=1$, giving
$X_G(X;q)=e_1(X)^2$. Even in this small case, evaluating a monomial
symmetric function requires subtraction; positivity of the generators
$E_k$ by itself does not settle the elementary coefficients.
\end{example}

\begin{lemma}[Color reversal and shape reciprocity]
\label{graph:reciprocity}
For every natural unit interval graph $G$, $X_G(X;q)=\chi_G(X;q)$.
Writing $c_\lambda(q)=[e_\lambda]X_G(X;q)$, one has
\begin{equation}
 c_\lambda(q)=q^M c_\lambda(q^{-1})
 \label{graph:reciprocity-eq}
\end{equation}
for every partition $\lambda$.
\end{lemma}

\begin{proof}
Use $r$ colors and send $f$ to $r+1-f$. This interchanges color
ascents and inversions and reverses the variable alphabet. Symmetry,
which follows from Proposition~\ref{graph:kernel} with
$\Gamma=G$, removes the variable reversal and gives
$X_G=\chi_G$. Since a proper coloring has
$\operatorname{asc}_G(f)+\operatorname{inv}_G(f)=M$, the same
argument gives $X_G(X;q)=q^M X_G(X;q^{-1})$. Take $r\ge n$, so
the elementary functions of degree $n$ are independent, and compare
their coefficients. Passing to the symmetric-function limit proves
Equation~\eqref{graph:reciprocity-eq}.
\end{proof}

It now suffices to prove positivity of
$[X_u]m_\lambda(Y)$ for a suitable ambient graph. The graph used in
Section~\ref{tri:section} has $D$ independent anchors and lies in an
interval of length $D-1$; independent vertices are spaced by at least
one, so its independence number is $D$.
Lemma~\ref{tri:embedding} places every original graph among its
vertices. The next two sections develop the positive wall sections
with which these monomial symmetric functions will be identified.
